\documentclass[11pt]{article}

% Change "review" to "final" to generate the final (sometimes called camera-ready) version.
% Change to "preprint" to generate a non-anonymous version with page numbers.
\usepackage[review]{acl}
\usepackage{booktabs}
\usepackage{multirow}
\usepackage{amsmath}

% Standard package includes
\usepackage{times}
\usepackage{latexsym}

% For proper rendering and hyphenation of words containing Latin characters (including in bib files)
\usepackage[T1]{fontenc}
% For Vietnamese characters
% \usepackage[T5]{fontenc}
% See https://www.latex-project.org/help/documentation/encguide.pdf for other character sets

% This assumes your files are encoded as UTF8
\usepackage[utf8]{inputenc}

% This is not strictly necessary, and may be commented out,
% but it will improve the layout of the manuscript,
% and will typically save some space.
\usepackage{microtype}

% This is also not strictly necessary, and may be commented out.
% However, it will improve the aesthetics of text in
% the typewriter font.
\usepackage{inconsolata}

%Including images in your LaTeX document requires adding
%additional package(s)
\usepackage{graphicx}

% If the title and author information does not fit in the area allocated, uncomment the following
%
%\setlength\titlebox{<dim>}
%
% and set <dim> to something 5cm or larger.

\title{Enigma@MultiClinAI: Multilingual Clinical Entity Recognition}

% Author information can be set in various styles:
% For several authors from the same institution:
% \author{Author 1 \and ... \and Author n \\
%         Address line \\ ... \\ Address line}
% if the names do not fit well on one line use
%         Author 1 \\ {\bf Author 2} \\ ... \\ {\bf Author n} \\
% For authors from different institutions:
% \author{Author 1 \\ Address line \\  ... \\ Address line
%         \And  ... \And
%         Author n \\ Address line \\ ... \\ Address line}
% To start a separate ``row'' of authors use \AND, as in
% \author{Author 1 \\ Address line \\  ... \\ Address line
%         \AND
%         Author 2 \\ Address line \\ ... \\ Address line \And
%         Author 3 \\ Address line \\ ... \\ Address line}

\author{First Author \\
  Affiliation / Address line 1 \\
  Affiliation / Address line 2 \\
  Affiliation / Address line 3 \\
  \texttt{email@domain} \\\And
  Second Author \\
  Affiliation / Address line 1 \\
  Affiliation / Address line 2 \\
  Affiliation / Address line 3 \\
  \texttt{email@domain} \\}

% Sylvia

% Order by number of submissions, f1 score
% Teodor
% Monika
% Vitosh
% Plamena
% Daniel

% Yulia
% Viktoria
% Koychev
% Boytcheva

%\author{
%  \textbf{First Author\textsuperscript{1}},
%  \textbf{Second Author\textsuperscript{1,2}},
%  \textbf{Third T. Author\textsuperscript{1}},
%  \textbf{Fourth Author\textsuperscript{1}},
%\\
%  \textbf{Fifth Author\textsuperscript{1,2}},
%  \textbf{Sixth Author\textsuperscript{1}},
%  \textbf{Seventh Author\textsuperscript{1}},
%  \textbf{Eighth Author \textsuperscript{1,2,3,4}},
%\\
%  \textbf{Ninth Author\textsuperscript{1}},
%  \textbf{Tenth Author\textsuperscript{1}},
%  \textbf{Eleventh E. Author\textsuperscript{1,2,3,4,5}},
%  \textbf{Twelfth Author\textsuperscript{1}},
%\\
%  \textbf{Thirteenth Author\textsuperscript{3}},
%  \textbf{Fourteenth F. Author\textsuperscript{2,4}},
%  \textbf{Fifteenth Author\textsuperscript{1}},
%  \textbf{Sixteenth Author\textsuperscript{1}},
%\\
%  \textbf{Seventeenth S. Author\textsuperscript{4,5}},
%  \textbf{Eighteenth Author\textsuperscript{3,4}},
%  \textbf{Nineteenth N. Author\textsuperscript{2,5}},
%  \textbf{Twentieth Author\textsuperscript{1}}
%\\
%\\
%  \textsuperscript{1}Affiliation 1,
%  \textsuperscript{2}Affiliation 2,
%  \textsuperscript{3}Affiliation 3,
%  \textsuperscript{4}Affiliation 4,
%  \textsuperscript{5}Affiliation 5
%\\
%  \small{
%    \textbf{Correspondence:} \href{mailto:email@domain}{email@domain}
%  }
%}

\begin{document}
\maketitle
\begin{abstract}

\end{abstract}

\section{Introduction}

\section{Related Work}
% Svetla

\section{Data}
% MultiClinNER dataset
% Sylvia
% Custom data used?
% Data augmentation

% [Teodor]: Adding a subsection which should probaly be shortened 
\subsection{Language - Czech}

The Czech training data derives from the DisTEMIST, SympTEMIST, and MedProcNER corpora
translated and annotation-projected to Czech~\cite{multiclinai2026}.
Documents are clinical case reports annotated in BRAT standoff format with character-offset
entity spans. Each entity type has a separate document set due to translation revision
differences, precluding multi-label training.

\begin{table}[h]
\centering
\small
\begin{tabular}{lrrrr}
\toprule
\textbf{Entity} & \multicolumn{2}{c}{\textbf{Documents}} & \multicolumn{2}{c}{\textbf{Entities}} \\
\cmidrule(lr){2-3}\cmidrule(lr){4-5}
 & \textbf{Train} & \textbf{Dev} & \textbf{Train} & \textbf{Dev} \\
\midrule
Disease   & 1,006 & 252 & 21,083 & 4,710 \\
Symptom   & 1,006 & 252 & 22,525 & 5,281 \\
Procedure & 1,006 & 252 & 22,438 & 5,063 \\
\bottomrule
\end{tabular}
\caption{Czech dataset statistics (entity counts from gold \texttt{.ann} files).}
\label{tab:cz-data}
\end{table}

\paragraph{Morphological synonym augmentation.}
Czech is morphologically rich: the same medical concept may appear in many inflected surface
forms (e.g., \textit{hypertenze} $\to$ \textit{hypertenz\'{i}}, \textit{hypertenzi}).
We build a Czech medical synonym dictionary (1,399 entries, 887 morphological families) and
generate augmented documents by replacing rare entity surface forms (frequency $\leq$5)
with morphological variants, yielding 2,161--2,271 new documents per entity type---more
than tripling the training set size.

\paragraph{LLM paraphrase augmentation.}
We additionally generated paraphrases using GPT-4.1-mini following the BioSynNER
methodology~\cite{multiclinai2026}, prompting the model to rewrite clinical sentences
while preserving entity spans. This provided a further 0.002--0.004 F1 gain on the
development set beyond morphological augmentation.

\paragraph{External data (negative result).}
We translated the NCBI Disease corpus (792 EN documents) to Czech via MarianMT with
marker-based entity alignment, but training on this silver data alone (F1=0.677) or
combined with morphological augmentation (F1=0.689) both underperformed morphological
augmentation alone (F1=0.710), indicating MT projection noise outweighs the coverage
benefit. DisTEMIST/SympTEMIST were excluded as data leakage (same source documents as
the competition training set).

% [Teodor]: End of subsection 

\section{Methods}
% Models used, architectures, training approach, hyperparameters, why/how were these models selected
% Links to huggingface to exact models used, cite paper if available

\subsection{Language - Italian}

\subsection{Language - Spanish}

\subsection{Language - English}

\subsection{Language - Romanian}

\subsection{Language - Dutch}

\subsection{Language - Czech}

\label{sec:methods-cz}
% [Teodor]: Again, the volume can be reduced
We train separate token classifiers for each of the three Czech entity types (disease,
symptom, procedure) using two pretrained transformer models:
\textbf{XLM-RoBERTa-base}~\cite{conneau-etal-2020-unsupervised} (\texttt{FacebookAI/xlm-roberta-base}),
a multilingual model covering 100 languages, and
\textbf{RobeCzech}~\cite{straka-strakova-2020-robeczech} (\texttt{ufal/robeczech-base}),
a Czech-specific RoBERTa model with an optimised Czech vocabulary.
Both are fine-tuned with a linear token classification head in a BIO tagging scheme;
we also evaluate an optional \textbf{Conditional Random Field (CRF)}
layer~\cite{lafferty2001conditional} that enforces valid BIO transitions via Viterbi
decoding, preventing invalid label sequences (e.g., \texttt{I-X} after \texttt{O}).

\paragraph{Training setup.}
Hyperparameters: learning rate $2{\times}10^{-5}$ (AdamW), weight decay 0.01, batch
size 16, max sequence length 512, early stopping patience 3 on development F1.
Final submission models are retrained on train+dev combined for 10 fixed epochs.
Rather than full documents, we split each document into sentences using punctuation
heuristics with a curated list of Czech abbreviations (e.g., \textit{mudr.},
\textit{např.}), retaining pure-\texttt{O} sentences to prevent recall bias (removing
them drops precision to 0.418). Entity-containing documents are oversampled once
(OS$\times$1) before merging with augmented data.

% [Teodor]: Not sure about what to do in general with this bug/fix
\paragraph{Tokeniser artifact fix.}
XLM-RoBERTa's SentencePiece byte-fallback emits overlapping tokens at identical character
offsets ($\sim$4 per Czech clinical document), causing contradictory B$\to$B training
signal and spurious short predicted spans.
During \textit{training}, overlapping tokens receive label $-100$ and are excluded from
the loss. During \textit{inference}, such tokens extend rather than split the current
span. As a \textit{post-filter}, spans of $\leq$1 character are discarded, and 2-character
spans are retained only if fully uppercase (valid abbreviations such as \textit{DM},
\textit{HT}).

\paragraph{Ensemble.}
We combine model checkpoints via span-level weighted voting (weights proportional to
development F1). For disease, four models are combined ($t{=}0.4$):
XLM-R+morph+oversample, RobeCzech+morph+paraphrase, RobeCzech+morph, and GPT-4.1-mini
(zero-shot). Despite low individual F1 (0.384), GPT-4.1-mini finds 224 entities missed
by all fine-tuned models (65.8\% complementarity), boosting ensemble recall from 0.719
to 0.730. Symptom and procedure use two-model ensembles
(XLM-R CRF + RobeCzech+oversample).


\subsection{Language - Swedish}

\section{Experiments and Results}
% Experiments with different models on dev/test set, discuss performance of different models and architectures
\subsection{Language - Italian}

\subsection{Language - Spanish}

\subsection{Language - English}

\subsection{Language - Romanian}

\subsection{Language - Dutch}

\subsection{Language - Czech}
\label{sec:results-cz}

% [Teodor]: Adding most things here, so we can filter easily later if/when needed. Derived from a "full paper" in czech_standalone.tex where most experiments done in the timespan of the cz competition are detaily listed. Also note that we should probaly remove the commenters names :) before submitting [Teodor]

\paragraph{Ablation (development set).}
Table~\ref{tab:cz-ablation} reports the development-set disease NER ablation;
symptom and procedure show consistent trends.
\textbf{Morphological augmentation} is the single most impactful intervention,
adding $\approx$4 F1 points over the non-augmented baseline (0.657 $\to$ 0.710--0.712).
\textbf{CRF decoding} adds a small consistent gain over softmax.
Despite being Czech-specific, RobeCzech slightly lags XLM-RoBERTa under morphological
augmentation alone (0.698 vs.\ 0.710), but benefits more from LLM paraphrases, reaching
0.714 as the best single model.
\textbf{Oversampling} (OS$\times$1) yields no additional gain beyond morphological
augmentation, though it helps without augmentation.

\begin{table}[h]
\centering
\small
\begin{tabular}{llllr}
\toprule
\textbf{Model} & \textbf{Aug.} & \textbf{Head} & \textbf{Extra} & \textbf{F1} \\
\midrule
XLM-R-base & None       & Softmax & {--}        & 0.657 \\
XLM-R-base & Morph      & Softmax & {--}        & 0.710 \\
XLM-R-base & Morph      & CRF     & {--}        & 0.712 \\
XLM-R-base & Morph      & Softmax & OS$\times$1 & 0.710 \\
XLM-R-base & Morph+Para & Softmax & {--}        & 0.704 \\
RobeCzech  & Morph      & Softmax & {--}        & 0.698 \\
RobeCzech  & Morph+Para & Softmax & {--}        & \textbf{0.714} \\
\bottomrule
\end{tabular}
\caption{Czech disease development F1 ablation. Morph = morphological synonym augmentation;
Para = LLM paraphrase; OS$\times$1 = entity-document oversampling.}
\label{tab:cz-ablation}
\end{table}

\paragraph{Ensemble (development set).}
Table~\ref{tab:cz-ensemble} compares ensemble strategies on the disease development set.
Weighted voting at $t{=}0.4$ achieves F1=0.738, outperforming the best single model by
2.4 points. Including GPT-4.1-mini increases recall (0.719 $\to$ 0.730) without
substantially hurting precision, since 65.8\% of its correct finds are entities missed
by all fine-tuned models.

\begin{table}[h]
\centering
\small
\begin{tabular}{llrrr}
\toprule
\textbf{System} & \textbf{Strategy} & \textbf{P} & \textbf{R} & \textbf{F1} \\
\midrule
\multirow{3}{*}{4-model}
  & Weighted $t{=}0.4$ & 0.747 & 0.730 & \textbf{0.738} \\
  & Threshold $K{=}2$  & 0.743 & 0.729 & 0.736 \\
  & Majority $K{=}3$   & 0.820 & 0.629 & 0.712 \\
3-model & Threshold $K{=}2$  & 0.753 & 0.719 & 0.736 \\
Best single & {--}             & 0.718 & 0.711 & 0.714 \\
\bottomrule
\end{tabular}
\caption{Disease development set ensemble results (4,710 gold entities).
4-model: XLM-R+morph+OS, RobeCzech+morph+para, RobeCzech+morph, GPT-4.1-mini.}
\label{tab:cz-ensemble}
\end{table}

\paragraph{Tokeniser fix impact.}
The SentencePiece byte-fallback fix brought substantial precision gains
(Table~\ref{tab:cz-bugfix}), confirming that overlapping token positions were generating
large numbers of spurious short entity spans. Disease precision improved from 0.565 to
0.708; symptom from 0.564 to 0.695.

\begin{table}[h]
\centering
\small
\begin{tabular}{llrrr}
\toprule
\textbf{Entity} & \textbf{System} & \textbf{P} & \textbf{R} & \textbf{F1} \\
\midrule
\multirow{2}{*}{Disease}
  & Ensemble (pre-fix)  & 0.565 & 0.674 & 0.615 \\
  & Ensemble (post-fix) & 0.708 & 0.675 & \textbf{0.691} \\
\midrule
\multirow{2}{*}{Procedure}
  & Ensemble (pre-fix)  & 0.617 & 0.685 & 0.649 \\
  & Ensemble (post-fix) & 0.634 & 0.692 & \textbf{0.662} \\
\midrule
\multirow{2}{*}{Symptom}
  & Ensemble (pre-fix)  & 0.564 & 0.629 & 0.595 \\
  & Ensemble (post-fix) & 0.695 & 0.582 & \textbf{0.634} \\
\bottomrule
\end{tabular}
\caption{Test set impact of the SentencePiece tokeniser artifact fix.}
\label{tab:cz-bugfix}
\end{table}

\paragraph{Test set results.}
Table~\ref{tab:cz-test} shows official test set results. The ensemble consistently
outperforms individual models. The largest gain from ensembling is for symptom (+0.038 F1),
where the individual model's lower recall is partially recovered by the complementary
RobeCzech model.

\begin{table}[h]
\centering
\small
\begin{tabular}{llrrr}
\toprule
\textbf{Entity} & \textbf{System} & \textbf{P} & \textbf{R} & \textbf{F1} \\
\midrule
\multirow{3}{*}{Disease}
  & Ensemble (4-model)  & \textbf{0.708} & \textbf{0.675} & \textbf{0.691} \\
  & XLM-R CRF+morph     & 0.651          & 0.638          & 0.645 \\
  & XLM-R+morph+OS      & 0.650          & 0.639          & 0.644 \\
\midrule
\multirow{2}{*}{Procedure}
  & Ensemble (2-model)  & 0.634          & \textbf{0.692} & \textbf{0.662} \\
  & XLM-R CRF+morph     & \textbf{0.668} & 0.655          & 0.662 \\
\midrule
\multirow{2}{*}{Symptom}
  & Ensemble (2-model)  & \textbf{0.695} & 0.582 & \textbf{0.634} \\
  & XLM-R CRF+morph     & 0.612          & 0.581 & 0.596 \\
\bottomrule
\end{tabular}
\caption{Czech MultiClinNER test set results (strict span matching, post-tokeniser-fix).}
\label{tab:cz-test}
\end{table}

\paragraph{Negative results.}
A CZ$\to$EN$\to$NER pipeline (MarianMT + Stanza i2b2 + awesome-align projection)
achieved only F1=0.107 on disease dev due to English NER over-prediction and span
fragmentation, and was not used.
Wikidata true-synonym augmentation (F1=0.693) slightly underperformed morphological-only
(F1=0.697), as concept-level substitution diverged from the clinical register.
XLM-RoBERTa-large yielded only +0.4\% F1 over the base model at twice the training cost
and was excluded.

\paragraph{Error analysis.}
GPT-4.1-mini uniquely recovers 224 disease entities missed by all fine-tuned models
(65.8\% complementarity), predominantly: rare multi-word terms (e.g.,
\textit{mezenchymální tumor}), medical abbreviations (e.g., \textit{PKU}), and
contextually implicit disease mentions.
Approximately 24\% of gold disease entities are missed by all models. Manual
inspection identifies very long multi-clause entity spans, highly inflected rare
forms with $\leq$2 training examples even after augmentation, negated or hypothetical
mentions (\textit{bez}, \textit{vyloučit}), and English--Czech code-mixed abbreviations
(\textit{COPD}, \textit{GERD}).
False positives are predominantly adjacent medical terms merged into a single span
and anatomical/test vocabulary misclassified as disease mentions.


\subsection{Language - Swedish}

\section{Conclusion}

% [Teodor]: If 1 thing should be said
For Czech, our key finding is that \textbf{morphological synonym augmentation} is the
most impactful single intervention (+4\% F1), generating morphologically varied training
documents for rare entity surface forms in a highly inflectional language.

\section*{Limitations}

\section*{Acknowledgments}

% Bibliography entries for the entire Anthology, followed by custom entries
%\bibliography{anthology,custom}
% Custom bibliography entries only
\bibliography{custom}

%\appendix

%\section{Example Appendix}
%\label{sec:appendix}


\end{document}
